%20_Theory.tex -- Rozdział 2. Podstawy teoretyczne
\tikzset{
  samolot/.pic={
    \fill[black!80] (-0.62,-0.06) -- (0.42,-0.06)
      .. controls (0.60,-0.06) and (0.62,0.06) .. (0.42,0.06) -- (-0.62,0.06) -- cycle;
    \fill[black!80] (-0.62,0.06) -- (-0.74,0.32) -- (-0.60,0.32) -- (-0.44,0.06) -- cycle;
    \fill[black!55] (0.12,-0.02) -- (-0.18,-0.30) -- (-0.06,-0.30) -- (0.26,-0.02) -- cycle;
    \fill[black!55] (-0.50,-0.01) -- (-0.66,-0.14) -- (-0.58,-0.14) -- (-0.40,-0.01) -- cycle;
  },
  satelita/.pic={
    \fill[black!75] (-0.09,-0.09) rectangle (0.09,0.09);
    \draw[black!75] (-0.13,0) -- (-0.09,0) (0.09,0) -- (0.13,0);
    \fill[akcent!60] (-0.42,-0.06) rectangle (-0.13,0.06);
    \fill[akcent!60] (0.13,-0.06) rectangle (0.42,0.06);
  },
  radar/.pic={
    \draw[black!80, line width=0.9pt] (0,0) -- (0,0.35);
    \fill[black!80] (-0.32,0.30) .. controls (-0.12,0.62) and (0.18,0.68) .. (0.36,0.62)
      -- (0.33,0.55) .. controls (0.15,0.58) and (-0.08,0.52) .. (-0.25,0.27) -- cycle;
    \draw[black!80, line width=0.9pt] (-0.3,0) -- (0.3,0);
  },
  odbiornik/.pic={
    \draw[black!80, line width=0.9pt] (0,0.25) -- (0,0.85);
    \fill[black!80] (0,0.85) circle (0.035);
    \fill[black!80] (-0.25,0) rectangle (0.25,0.25);
  },
}


% ---------------------------------------------------------------------
\chapter{Podstawy teoretyczne}
\label{chap:teoria}
% =====================================================================


W~rozdziale przedstawiono podstawy teoretyczne zagadnień wykorzystanych w~dalszej części pracy. Omówiono system ADS-B i~budowę ramek Mode~S, sposób kodowania położenia w~wiadomościach ADS-B, odbiór sygnałów za pomocą radia definiowanego programowo, wyznaczanie pozycji i~czasu przez odbiornik GNSS oraz środowisko ROS~2, w~którym udostępniane są dane. Rozdział zamyka przegląd istniejących rozwiązań.


% ---------------------------------------------------------------------
\section{Standard ADS-B i~ramki Mode~S}
% =====================================================================
\label{sec:adsb}


Początkiem dozorowania ruchu lotniczego było wykorzystanie radaru pierwotnego (ang.~\textit{primary surveillance radar}, PSR). Radar ten wysyła krótkie impulsy radiowe i~odbiera ich odbicia od statku powietrznego. Na podstawie różnicy czasu między wysłaniem impulsu a~odebraniem echa radar wyznacza odległość do celu, a~azymut z~kąta obrotu anteny. Rozwiązanie to nie dostarcza jednak informacji o~tożsamości statku powietrznego ani o~jego wysokości~\cite{sun2021}.

Brakujące informacje dostarcza radar wtórny (ang.~\textit{secondary surveillance radar}, SSR), który wysyła zapytania na częstotliwości \SI{1030}{\mega\hertz}. Transponder statku powietrznego odpowiada na nie na częstotliwości \SI{1090}{\mega\hertz}. W~pierwszych trybach pracy, oznaczonych literami A i~C, odpowiedź transpondera zawiera w~trybie~A czterocyfrowy kod identyfikacyjny w~zapisie ósemkowym (ang.~\textit{squawk code}), a~w~trybie~C wysokość barometryczną. O~ile rozwiązanie to sprawdzało się przy niewielkim natężeniu ruchu, o~tyle wraz z~rozwojem lotnictwa ujawniły się jego ograniczenia. Gdy kilka statków powietrznych znajduje się w~tym samym kierunku względem radaru, ich odpowiedzi nakładają się na siebie, co uniemożliwia poprawny odczyt. Dodatkowo liczba dostępnych kodów, wynosząca \num{4096} ($8^4$), okazała się niewystarczająca dla obecnego ruchu lotniczego~\cite{sun2021}.

Część z~tych problemów rozwiązano, wprowadzając tryb Mode~S (ang.~\textit{Mode Select}), opracowany w~latach siedemdziesiątych XX~wieku w~Lincoln Laboratory Massachusetts Institute of Technology i~wdrażany w~kontroli ruchu lotniczego od lat dziewięćdziesiątych. Najważniejszą cechą wyróżniającą ten tryb jest wysyłanie selektywnych zapytań, kierowanych do wybranego statku powietrznego, co ograniczyło problem nakładania się odpowiedzi. Dodatkowo zamiast jednego z~\num{4096} kodów każdy transponder Mode~S ma unikalny 24-bitowy adres, nazywany adresem ICAO, co pozwala rozróżnić do \num{16777216} statków powietrznych. Odpowiedzi Mode~S mają postać ramek o~długości 56 lub 112~bitów, przy czym pierwsze 5~bitów (ang.~\textit{downlink format}, DF) jednoznacznie określa strukturę i~interpretację pozostałej części ramki~\cite{sun2021}.

W~trybie Mode~S przewidziano również tzw. rozszerzony squitter (ang.~\textit{extended squitter}). Jest to długa, 112-bitowa ramka wysyłana przez transponder samoczynnie, bez uprzedniego zapytania ze strony radaru~\cite{sun2021}. Ramki tego typu przenoszą wiadomości systemu ADS-B (ang.~\textit{Automatic Dependent Surveillance--Broadcast}). W~systemie tym statek powietrzny wyznacza własną pozycję za pomocą odbiornika nawigacji satelitarnej (GNSS), a~następnie regularnie rozgłasza ją drogą radiową wraz z~prędkością i~innymi parametrami lotu. Wiadomości ADS-B nadaje ten sam transponder Mode~S i~na tej samej częstotliwości \SI{1090}{\mega\hertz}, na której wysyłane są odpowiedzi dla radaru wtórnego. Oba systemy różnią się jednak zasadą dozorowania: w~radarze wtórnym położenie wyznacza stacja naziemna, a~w~systemie ADS-B statek powietrzny. Różnicę tę przedstawiono na rysunku~\ref{rys:ssr_adsb}.
 
\begin{figure}[htbp]
  \centering
  \begin{tikzpicture}[font=\footnotesize, >=Stealth]
    \node[anchor=west] at (0,4.35) {(a) radar wtórny};
    \draw[black!50] (0,0) -- (6.6,0);
    \pic at (1.1,0) {radar};
    \pic[scale=1.1] at (5.3,3.25) {samolot};
    \draw[->, akcent, thick] (1.50,0.80) -- node[sloped, above] {zapytanie, \SI{1030}{\mega\hertz}} (4.55,3.05);
    \draw[->, black!80, thick] (4.65,2.70) -- node[sloped, below] {odpowiedź, \SI{1090}{\mega\hertz}} (1.65,0.55);
    \node[align=center, anchor=north] at (1.1,-0.08) {radar mierzy odległość\\ i~azymut};
    \node[align=center, anchor=south] at (5.3,3.55) {transponder podaje\\ kod lub wysokość};
    \begin{scope}[xshift=7.7cm]
      \node[anchor=west] at (0,4.35) {(b) ADS-B};
      \draw[black!50] (0,0) -- (7.2,0);
      \pic at (0.9,3.55) {satelita};
      \pic at (3.6,3.90) {satelita};
      \pic at (6.3,3.55) {satelita};
      \node[black!70] at (5.45,4.25) {satelity GNSS};
      \draw[->, black!50, dashed] (1.10,3.40) -- (3.15,2.55);
      \draw[->, black!50, dashed] (3.60,3.75) -- (3.60,2.70);
      \draw[->, black!50, dashed] (6.10,3.40) -- (4.05,2.55);
      \pic[scale=1.1] at (3.6,2.35) {samolot};
      \pic at (0.9,0) {odbiornik};
      \node[anchor=north] at (0.9,-0.08) {stacja bazowa};
      \pic[scale=0.8] at (6.3,1.25) {samolot};
      \node[anchor=north, align=center] at (6.3,0.95) {inny statek\\ powietrzny};
      \draw[->, akcent, thick] (3.10,1.95) -- (1.20,0.98);
      \draw[->, akcent, thick] (4.10,1.95) -- (5.65,1.35);
      \node[align=center, akcent] at (3.60,0.75) {rozgłaszanie, \SI{1090}{\mega\hertz}:\\ pozycja, prędkość,\\ identyfikacja};
    \end{scope}
  \end{tikzpicture}
  \caption{Zasada dozorowania: (a) radar wtórny -- położenie wyznacza stacja naziemna, (b) system ADS-B -- położenie wyznacza i~rozgłasza statek powietrzny; opracowanie własne na podstawie~\cite{sun2021}}
  \label{rys:ssr_adsb}
\end{figure}
W~przestrzeni powietrznej Unii Europejskiej nadajnik ADS-B pracujący na częstotliwości \SI{1090}{\mega\hertz} w~trybie rozszerzonego squittera stanowi obowiązkowe wyposażenie statków powietrznych o~maksymalnej masie startowej powyżej \SI{5700}{kg} lub o~maksymalnej prędkości przelotowej przekraczającej \SI{250}{kt}~\cite{ue2023_1770}. W~efekcie wiadomości ADS-B nadaje obecnie większość statków powietrznych~\cite{sun2021}. Rozważano również inne sposoby realizacji ADS-B, m.in. łącze UAT (ang.~\textit{Universal Access Transceiver}). Najszerzej przyjęła się jednak realizacja oparta na rozszerzonym squitterze Mode~S, a~łącze UAT wdrożono jedynie częściowo, np. w~Stanach Zjednoczonych~\cite{sun2021}. W~dalszej części pracy pojęcie ADS-B odnosi się wyłącznie do realizacji opartej na rozszerzonym squitterze.

Ramka ADS-B ma długość 112~bitów i~jest podzielona na pięć pól, które przedstawiono na rysunku~\ref{rys:ramka_df17}~\cite{sun2021}.
 
\begin{figure}[htbp]
  \centering
  \begin{tikzpicture}[font=\footnotesize, x=0.9cm]
    \def\h{0.8}
    \foreach \a/\b/\n/\r in {0/1.1/DF/1--5, 1.1/2.1/CA/6--8, 2.1/5.6/ICAO/9--32,
                             5.6/6.8/TC/33--37, 6.8/11.6/ME (dane)/38--88, 11.6/15/PI/89--112} {
      \draw[thick, fill=akcent!8] (\a,0) rectangle (\b,\h);
      \node at ({(\a+\b)/2},{\h/2}) {\textbf{\n}};
      \node[above] at ({(\a+\b)/2},\h) {\r};
    }
    \fill[akcent!20] (5.6,0) rectangle (6.8,\h);
    \node at (6.2,{\h/2}) {\textbf{TC}};
    \draw[thick] (5.6,0) rectangle (11.6,\h);
    \draw[decorate, decoration={brace, mirror, amplitude=4pt}, black!70]
      (5.62,-0.08) -- (11.58,-0.08) node[midway, below=4pt] {pole ME (56 bitów)};
    \node[anchor=east, black!70] at (-0.15,{\h+0.22}) {bity:};
    \node[anchor=east, black!70] at (-0.15,-1.0) {hex:};
    \node[anchor=east, black!70] at (-0.15,-1.75) {znaczenie:};
    \node at (1.05,-1.0) {\texttt{8D}};
    \node at (3.85,-1.0) {\texttt{461E16}};
    \node at (8.6,-1.0)  {\texttt{231893B61E0820}};
    \node at (13.3,-1.0) {\texttt{3120B1}};
    \node[align=center] at (1.05,-1.75) {DF\,=\,17\\ CA\,=\,5};
    \node[align=center] at (3.85,-1.75) {adres ICAO\\ nadawcy};
    \node[align=center] at (8.6,-1.75)  {TC\,=\,4: identyfikacja,\\ znak wywoławczy \texttt{FIN6G}};
    \node[align=center] at (13.3,-1.75) {suma kontrolna\\ poprawna};
  \end{tikzpicture}
  \caption{Pola ramki ADS-B na przykładzie ramki \texttt{8D461E16231893B61E08203120B1} zarejestrowanej przez stację bazową; pierwsze 5~bitów pola ME stanowi kod typu TC; szerokości pól nie są proporcjonalne do liczby bitów; opracowanie własne na podstawie~\cite{sun2021}}
  \label{rys:ramka_df17}
\end{figure}
 
\newpage
Znaczenie pól ramki i przykładowe wartości zestawiono w~tabeli~\ref{tab:pola_ramki}.
\begin{table}[H]
  \centering
  \caption{Pola ramki ADS-B i interpretacja ich wartości; opracowanie własne na podstawie~\cite{sun2021}}
  \label{tab:pola_ramki}
  \small
  \begin{tabular}{@{}lcp{10.2cm}@{}}
    \toprule
    Pole & Bity & Znaczenie i przykładowe wartości \\
    \midrule
    DF & 1--5 & format ramki: 17 -- wiadomość ADS-B nadana przez transponder Mode~S, 18 -- wiadomość ADS-B z~urządzenia niebędącego transponderem, 11 -- krótki squitter, 4, 5, 20, 21 -- odpowiedzi na zapytania radaru wtórnego \\
    CA & 6--8 & możliwości transpondera: 0 -- poziom~1 (tylko dozorowanie), 4 -- poziom~2 lub wyższy, statek powietrzny na ziemi, 5 -- poziom~2 lub wyższy, statek powietrzny w~powietrzu, 6 -- poziom~2 lub wyższy, na ziemi lub w~powietrzu \\
    ICAO & 9--32 & 24-bitowy adres statku powietrznego; adresy są przydzielane państwom blokami, dlatego początek adresu wskazuje państwo rejestracji \\
    ME & 33--88 & 56 bitów danych; pierwsze 5~bitów to kod typu TC, który określa, jak interpretować pozostałe 51~bitów (położenie, prędkość, identyfikacja i~inne) \\
    PI & 89--112 & 24-bitowa suma kontrolna CRC obliczona z~pierwszych 88~bitów ramki \\
    \bottomrule
  \end{tabular}
\end{table}
W~przykładzie z~rysunku~\ref{rys:ramka_df17} wartość DF\,=\,17 oznacza wiadomość ADS-B. Wartość CA\,=\,5 wskazuje na transponder poziomu~2 lub wyższego, a~ponadto informuje, że statek powietrzny znajduje się w~powietrzu~\cite{sun2021}.

Pole ME ramki z~rysunku~\ref{rys:ramka_df17} zapisano szesnastkowo jako \texttt{231893B61E0820}. Pierwszy bajt, \texttt{23}, ma postać binarną \texttt{00100011}. Jego pierwsze pięć bitów (\texttt{00100}) daje kod typu TC\,=\,4, czyli wiadomość identyfikacyjną. Kolejne trzy bity (\texttt{011}) zawierają kategorię statku powietrznego. Według~\cite{sun2021} wartość 3 dla TC\,=\,4 oznacza średni statek powietrzny kategorii~2, czyli o~masie od 34 do 136~t. Pozostałe 48~bitów dzieli się na osiem bloków po 6~bitów, z~których każdy koduje jeden znak znaku wywoławczego. Wartości liczbowe bloków odpowiadają znakom zgodnie z~tablicą podaną  w~książce~\cite{sun2021}: liczby od 1 do 26 oznaczają litery A--Z, liczba 32 -- spację, a~liczby od 48 do 57 -- cyfry 0--9. Przebieg dekodowania przedstawiono w~tabeli~\ref{tab:znak_wywolawczy}. Wynikiem jest znak wywoławczy FIN6G, dopełniony trzema spacjami do ośmiu znaków.

\begin{table}[htbp]
  \centering
  \caption{Dekodowanie znaku wywoławczego z~pola ME ramki z~rysunku~\ref{rys:ramka_df17}; symbol \textvisiblespace{} oznacza spację; opracowanie własne na podstawie~\cite{sun2021}}
  \label{tab:znak_wywolawczy}
  \small
  \begin{tabular}{@{}lcccccccc@{}}
    \toprule
    Numer bloku & 1 & 2 & 3 & 4 & 5 & 6 & 7 & 8 \\
    \midrule
    Bity    & \texttt{000110} & \texttt{001001} & \texttt{001110} & \texttt{110110} & \texttt{000111} & \texttt{100000} & \texttt{100000} & \texttt{100000} \\
    Wartość & 6 & 9 & 14 & 54 & 7 & 32 & 32 & 32 \\
    Znak    & F & I & N & 6 & G & \textvisiblespace & \textvisiblespace & \textvisiblespace \\
    \bottomrule
  \end{tabular}
\end{table}
\newpage
Ostatnim elementem ramki jest 24-bitowe pole PI, zawierające cykliczną sumę kontrolną CRC (ang.~\textit{cyclic redundancy check}), obliczaną na podstawie pierwszych 88~bitów ramki~\cite{dump1090}. Odbiornik wyznacza własną wartość sumy dla odebranych danych i~porównuje ją z~wartością zapisaną w~polu PI. Niezgodność wykazana podczas weryfikacji świadczy o~uszkodzeniu ramki w~trakcie transmisji radiowej. Wykorzystany algorytm CRC umożliwia nie tylko wykrywanie przekłamań, ale również korekcję pojedynczych błędów bitowych~\cite{sun2021}.

Przy wykorzystaniu wiadomości ADS-B jako danych referencyjnych należy uwzględnić ograniczenia systemu. Położenie podaje sam statek powietrzny, a~protokół nie przewiduje uwierzytelnienia nadawcy ani szyfrowania wiadomości. System jest przez to podatny na celowe zakłócanie oraz nadawanie fałszywych wiadomości (ang.~\textit{spoofing})~\cite{strohmeier2015security}. Kanał \SI{1090}{\mega\hertz} jest współdzielony przez transpondery wszystkich statków powietrznych, które nadają zarówno wiadomości ADS-B, jak i~odpowiedzi na zapytania radarów wtórnych. Przy dużym natężeniu ruchu część ramek nakłada się więc na siebie i~zostaje utracona~\cite{strohmeier2014realities}. Adres ICAO może ponadto zostać przeprogramowany, co obserwowano m.in. w~przypadku statków powietrznych wojskowych~\cite{sun2021}. Wymagania wynikające z~tych ograniczeń sformułowano w~podrozdziale~\ref{sec:wymagania_referencja}.


% ---------------------------------------------------------------------
\section{Kodowanie położenia CPR}\label{sec:cpr}
% =====================================================================


W~wiadomości o~położeniu na każdą ze współrzędnych geograficznych przeznaczono 17~bitów~\cite{sun2021}. Bezpośredni zapis długości geograficznej z~zakresu \ang{360} na 17~bitach dawałby rozdzielczość około \ang{0.0027}, czyli na równiku około \SI{300}{m}. Rozdzielczość ta jest wielokrotnie gorsza od dokładności odbiorników GNSS. Z~tego powodu stosuje się kodowanie CPR (ang.~\textit{compact position reporting}), w~którym przesyła się położenie statku powietrznego wewnątrz strefy o~ograniczonych wymiarach. Dla strefy o~rozciągłości \ang{6} szerokości geograficznej ta sama liczba bitów zapewnia rozdzielczość około \SI{5}{m}. Ceną za dokładność jest niejednoznaczność, ponieważ pojedyncza wiadomość nie wskazuje, w~której strefie znajduje się statek powietrzny~\cite{sun2021}.

Statek powietrzny nadaje na przemian dwa rodzaje wiadomości o~położeniu: parzyste i~nieparzyste. Różnią się one rozmiarem stref szerokości geograficznej, który wynosi~\cite{sun2021}
\begin{equation}
  \Delta\varphi_i = \frac{\ang{360}}{4N_{\mathrm{Z}} - i},
  \label{eq:dlat}
\end{equation}
\begin{gdzie}
  \item $\Delta\varphi_i$ -- rozciągłość strefy szerokości geograficznej, w~stopniach,
  \item $N_{\mathrm{Z}} = 15$ -- liczba stref szerokości geograficznej między równikiem a~biegunem,
  \item $i$ -- rodzaj wiadomości: $i = 0$ dla wiadomości parzystej, $i = 1$ dla nieparzystej.
\end{gdzie}
Strefy wiadomości parzystych mają więc rozciągłość \ang{6}, a~nieparzystych -- około \ang{6.1}. Liczba stref długości geograficznej, oznaczana $\operatorname{NL}(\varphi)$, zależy od szerokości geograficznej $\varphi$. Wynosi ona 59 przy równiku i~maleje do 1 w~pobliżu biegunów~\cite{sun2021}.
Wyraża ją wzór~\cite{sun2021}
\begin{equation}
  \operatorname{NL}(\varphi) = \left\lfloor
    \frac{2\pi}{\arccos\left(1 - \dfrac{1 - \cos\dfrac{\pi}{2N_{\mathrm{Z}}}}{\cos^{2}\varphi}\right)}
  \right\rfloor,
  \label{eq:nl}
\end{equation}
w~którym symbol $\lfloor\cdot\rfloor$ oznacza zaokrąglenie w~dół do liczby całkowitej. Przebieg funkcji~\eqref{eq:nl} przedstawiono na rysunku~\ref{rys:nl}. Dla szerokości geograficznej Krakowa $\operatorname{NL} = 38$, co daje strefy o~rozciągłości około \ang{9.5} długości geograficznej.

\begin{figure}[htbp]
  \centering
  \begin{tikzpicture}
    \begin{axis}[
      width=13cm, height=6.2cm,
      xmin=0, xmax=90, ymin=0, ymax=62,
      xtick={0,10,...,90}, ytick={0,10,...,60},
      xlabel={$\varphi$ [\si{\degree}]}, ylabel={$\operatorname{NL}(\varphi)$ [--]},
      grid=major, grid style={black!12},
      tick label style={font=\footnotesize}, label style={font=\footnotesize},
      /pgf/number format/use comma,
    ]
      \addplot[const plot mark left, thick, akcent] coordinates {
        (0,59) (10.4705,58) (14.8282,57) (18.1863,56) (21.0294,55) (23.545,54) (25.8292,53)
        (27.939,52) (29.9114,51) (31.7721,50) (33.5399,49) (35.229,48) (36.8503,47) (38.4124,46)
        (39.9226,45) (41.3865,44) (42.8091,43) (44.1945,42) (45.5463,41) (46.8673,40) (48.1604,39)
        (49.4278,38) (50.6715,37) (51.8934,36) (53.0952,35) (54.2782,34) (55.4438,33) (56.5932,32)
        (57.7275,31) (58.8476,30) (59.9546,29) (61.0492,28) (62.1322,27) (63.2043,26) (64.2662,25)
        (65.3185,24) (66.3617,23) (67.3965,22) (68.4232,21) (69.4424,20) (70.4545,19) (71.4599,18)
        (72.4588,17) (73.4518,16) (74.4389,15) (75.4206,14) (76.3968,13) (77.3679,12) (78.3337,11)
        (79.2943,10) (80.2492,9) (81.198,8) (82.1396,7) (83.072,6) (83.9917,5) (84.8917,4)
        (85.7554,3) (86.5354,2) (87,1) (90,1)
      };
      \draw[dashed, black!70] (axis cs:50.06,0) -- (axis cs:50.06,38);
      \node[anchor=south west, font=\footnotesize] at (axis cs:51,39) {Kraków: $\varphi \approx \ang{50.1}$, $\operatorname{NL} = 38$};
    \end{axis}
  \end{tikzpicture}
  \caption{Liczba stref długości geograficznej w~funkcji szerokości geograficznej; opracowanie własne na podstawie wzoru~\eqref{eq:nl}}
  \label{rys:nl}
\end{figure}

W~wiadomości przesyła się położenie statku powietrznego wewnątrz strefy, wyrażone ułamkiem jej rozmiaru~\cite{sun2021}:
\begin{equation}
  \varphi_{\mathrm{cpr},i} = \frac{Y_i}{2^{17}},
  \label{eq:cpr}
\end{equation}
\begin{gdzie}
  \item $\varphi_{\mathrm{cpr},i}$ -- położenie wewnątrz strefy szerokości geograficznej, z~przedziału $[0, 1)$,
  \item $Y_i$ -- 17-bitowa liczba całkowita odczytana z~wiadomości.
\end{gdzie}
Długość geograficzną koduje się analogicznie. Pojedyncza wiadomość wskazuje więc tylko zbiór możliwych położeń, po jednym w~każdej strefie. Do wyboru właściwego z~nich stosuje się dekodowanie globalne albo lokalne, których ideę przedstawiono w~uproszczonej, jednowymiarowej postaci na rysunku~\ref{rys:cpr_idea}.

\begin{figure}[htbp]
  \centering
  \begin{tikzpicture}[font=\footnotesize, x=0.38cm, y=1cm, >=Stealth]
    \node[anchor=west] at (-6.0,3.55) {(a) dekodowanie globalne};
    \node[anchor=east] at (-0.4,2.8) {wiadomość parzysta};
    \draw[black!60] (0,2.8) -- (30,2.8);
    \foreach \x in {0,6,...,30} \draw[black!60] (\x,2.65) -- (\x,2.95);
    \foreach \x in {5,11,17,23,29} \fill[akcent] (\x,2.8) circle (2.3pt);
    \node[anchor=east] at (-0.4,2.0) {wiadomość nieparzysta};
    \draw[black!60] (0,2.0) -- (30,2.0);
    \foreach \x in {0,5,...,30} \draw[black!60] (\x,1.85) -- (\x,2.15);
    \foreach \x in {2,7,12,17,22,27} \fill[black!70] (\x,2.0) circle (2.3pt);
    \draw[dashed, thick] (17,1.7) -- (17,3.2);
    \node[anchor=south] at (17,3.2) {jedyny wspólny kandydat};
    \node[anchor=west] at (-6.0,1.15) {(b) dekodowanie lokalne};
    \fill[akcent!12] (12,0.2) rectangle (18,0.8);
    \node[anchor=east] at (-0.4,0.5) {wiadomość parzysta};
    \draw[black!60] (0,0.5) -- (30,0.5);
    \foreach \x in {0,6,...,30} \draw[black!60] (\x,0.35) -- (\x,0.65);
    \foreach \x in {5,11,23,29} \fill[akcent!35] (\x,0.5) circle (2.3pt);
    \fill[akcent] (17,0.5) circle (2.8pt);
    \draw[->, thick] (15,-0.25) -- (15,0.33);
    \node[anchor=north] at (15,-0.25) {punkt odniesienia};
    \node[anchor=south, black!70] at (15,0.8) {$\pm$ pół strefy};
    \draw[->, black!70] (0,-1.05) -- (31.2,-1.05);
    \node[anchor=north] at (15,-1.5) {położenie [j.u.]};
    \foreach \x in {0,5,...,30} {
      \draw[black!70] (\x,-1.05) -- (\x,-1.12);
      \node[anchor=north] at (\x,-1.12) {\x};
    }
  \end{tikzpicture}
  \caption{Idea dekodowania położenia CPR w~jednym wymiarze, dla stref o~rozmiarach 6 i~5~jednostek umownych (j.u.): (a) dekodowanie globalne -- położenie wyznacza jedyny wspólny kandydat wiadomości parzystej i~nieparzystej, (b) dekodowanie lokalne -- położenie wyznacza kandydat leżący bliżej niż pół strefy od punktu odniesienia; kreski oznaczają granice stref, a~kropki -- możliwe położenia; opracowanie własne na podstawie~\cite{sun2021}}
  \label{rys:cpr_idea}
\end{figure}

W~dekodowaniu globalnym wykorzystuje się parę wiadomości: parzystą i~nieparzystą. Strefy obu rodzajów wiadomości mają różne rozmiary, dlatego tylko jedno z~możliwych położeń wskazywanych przez wiadomość parzystą pokrywa się z~położeniem wskazywanym przez wiadomość nieparzystą (rysunek~\ref{rys:cpr_idea}a). Obie wiadomości powinny zostać odebrane w~krótkim odstępie czasu, w~praktyce nie dłuższym niż około \SI{10}{s}. Statek powietrzny nie może przy tym przekroczyć granicy stref, dla której zmienia się wartość $\operatorname{NL}$~\cite{sun2021}. W~dekodowaniu lokalnym wystarcza jedna wiadomość, jeżeli znane jest położenie odniesienia leżące bliżej niż pół strefy od rzeczywistego położenia (rysunek~\ref{rys:cpr_idea}b). Dla wiadomości nadawanych w~powietrzu warunek ten oznacza odległość mniejszą niż \SI{180}{NM}~\cite{sun2021}, co odpowiada połowie strefy szerokości geograficznej: $\ang{3} \cdot \SI[per-mode=symbol]{111.2}{\kilo\metre\per\degree} \approx \SI{333}{km}$. Punktem odniesienia jest zwykle wcześniej wyznaczone położenie statku powietrznego~\cite{sun2021}. W~przypadku odbiornika naziemnego może nim być również położenie odbiornika, o~ile statek powietrzny znajduje się w~odległości mniejszej niż \SI{180}{NM}.

W~dekodowaniu globalnym oblicza się najpierw indeks strefy szerokości geograficznej, a~następnie szerokość geograficzną dla każdej z~wiadomości~\cite{sun2021}:
\begin{equation}
  j = \left\lfloor 59\,\varphi_{\mathrm{cpr},0} - 60\,\varphi_{\mathrm{cpr},1} + \frac{1}{2} \right\rfloor,
  \qquad
  \varphi_i = \Delta\varphi_i \bigl( \operatorname{mod}(j,\, 60 - i) + \varphi_{\mathrm{cpr},i} \bigr),
  \label{eq:globalne}
\end{equation}
\begin{gdzie}
  \item $j$ -- indeks strefy szerokości geograficznej,
  \item $\varphi_i$ -- szerokość geograficzna wyznaczona z~wiadomości rodzaju $i$, w~stopniach,
  \item $\operatorname{mod}(a, b)$ -- reszta z~dzielenia $a$ przez $b$, nieujemna dla $b > 0$.
\end{gdzie}
Wartości nie mniejsze niż \ang{270} odpowiadają półkuli południowej i~pomniejsza się je o~\ang{360}. Jako wynik przyjmuje się szerokość geograficzną obliczoną z~nowszej wiadomości pary~\cite{sun2021}. W~dekodowaniu lokalnym indeks strefy wyznacza się z~położenia odniesienia $\varphi_{\mathrm{ref}}$~\cite{sun2021}:
\begin{equation}
  j = \left\lfloor \frac{\varphi_{\mathrm{ref}}}{\Delta\varphi_i} \right\rfloor
    + \left\lfloor \frac{\operatorname{mod}(\varphi_{\mathrm{ref}},\, \Delta\varphi_i)}{\Delta\varphi_i}
    - \varphi_{\mathrm{cpr},i} + \frac{1}{2} \right\rfloor,
  \qquad
  \varphi = \Delta\varphi_i \left( j + \varphi_{\mathrm{cpr},i} \right).
  \label{eq:lokalne}
\end{equation}
Długość geograficzną oblicza się w~obu metodach analogicznie, przy czym rozmiar strefy wynosi $\ang{360} / \max\bigl(\operatorname{NL}(\varphi) - i,\, 1\bigr)$ i~zależy od wyznaczonej wcześniej szerokości geograficznej~\cite{sun2021}.

Przykładowo niech stacja bazowa znajduje się na szerokości geograficznej $\varphi_{\mathrm{ref}} = \ang{50.0}$, a~statek powietrzny na szerokości \ang{50.3169}, jak w~przykładzie z~podrozdziału~\ref{sec:model}. Wiadomość parzysta zawiera wówczas wartość $Y_0 = \num{50613}$, czyli $\varphi_{\mathrm{cpr},0} \approx \num{0.38615}$. Ze wzoru~\eqref{eq:lokalne} otrzymuje się $j = 8 + \lfloor \num{0.33333} - \num{0.38615} + \num{0.5} \rfloor = 8$ oraz $\varphi = \ang{6} \cdot (8 + \num{0.38615}) \approx \ang{50.31688}$. Wynik różni się od rzeczywistej szerokości geograficznej o~około \SI{2}{m}, co wynika z~rozdzielczości kodowania.

Algorytm CPR nie gwarantuje poprawności każdego wyniku i~w~rzadkich przypadkach prowadzi do błędnego położenia. Z~tego powodu zaleca się sprawdzanie wiarygodności wyznaczonych położeń, np. przez porównanie ich z~zasięgiem odbiornika lub z~położeniami obliczonymi z~innych par wiadomości~\cite{sun2021}.


%---------------------------------------------------------------------------
\section{Odbiór sygnałów z~wykorzystaniem SDR}
% =====================================================================


\label{sec:sdr}

Odpowiedzi Mode~S, w~tym ramki ADS-B, są przesyłane z~szybkością \SI{1}{Mbit/s}~\cite{sun2021}. Każda ramka rozpoczyna się preambułą o~czasie trwania \SI{8}{\micro\second}, po której następuje blok danych o~długości 56 lub \SI{112}{\micro\second}. Dane zakodowano za pomocą modulacji położenia impulsu (ang.~\textit{pulse position modulation}, PPM). Bit o~wartości 1 odpowiada impulsowi w~pierwszej połowie mikrosekundy, a~bit o~wartości 0 impulsowi w~drugiej połowie~\cite{sun2021}. Na rysunku~\ref{rys:ppm} przedstawiono początek ramki z~rysunku~\ref{rys:ramka_df17}.

\begin{figure}[htbp]
  \centering
  \begin{tikzpicture}[font=\footnotesize, x=0.8cm, y=1cm]
    \draw[black!60] (0,0) -- (16.3,0);
    \foreach \t in {0,...,16} \draw[black!60] (\t,0) -- (\t,-0.08);
    \foreach \t in {0,2,...,16} \node[anchor=north] at (\t,-0.08) {\t};
    \node[anchor=north west] at (16.35,-0.08) {$t$ [\si{\micro\second}]};
    \foreach \t in {0,1,3.5,4.5} \fill[akcent] (\t,0) rectangle ++(0.5,0.8);
    \foreach \k/\b in {0/1,1/0,2/0,3/0,4/1,5/1,6/0,7/1} {
      \draw[black!25, dashed] ({8+\k},0) -- ({8+\k},1.05);
      \ifnum\b=1
        \fill[black!75] ({8+\k},0) rectangle ++(0.5,0.8);
      \else
        \fill[black!75] ({8.5+\k},0) rectangle ++(0.5,0.8);
      \fi
      \node at ({8.5+\k},1.25) {\b};
    }
    \draw[black!25, dashed] (16,0) -- (16,1.05);
    \draw[decorate, decoration={brace, amplitude=5pt}] (0,1.55) -- (8,1.55)
      node[midway, above=5pt] {preambuła (\SI{8}{\micro\second})};
    \draw[decorate, decoration={brace, amplitude=5pt}] (8,1.55) -- (16,1.55)
      node[midway, above=5pt] {pierwsze 8~bitów danych: \texttt{8D} = \texttt{10001101}};
  \end{tikzpicture}
  \caption{Preambuła i~pierwsze bity ramki z~rysunku~\ref{rys:ramka_df17} zakodowane metodą PPM; opracowanie własne na podstawie~\cite{sun2021}}
  \label{rys:ppm}
\end{figure}

W~radiu definiowanym programowo (ang.~\textit{software-defined radio}, SDR) większość operacji przetwarzania sygnału, w~tym demodulację, wykonuje oprogramowanie. Część sprzętowa ogranicza się do odbioru sygnału, przeniesienia go na niższą częstotliwość i~próbkowania. Użytkownik wybiera programowo m.in. częstotliwość środkową odbiornika i~częstotliwość próbkowania~\cite{sun2021}. Najpopularniejszymi tanimi odbiornikami SDR są urządzenia RTL-SDR, zbudowane na układzie Realtek RTL2832U i~pierwotnie przeznaczone do odbioru telewizji cyfrowej DVB-T~\cite{sun2021,rtlsdrv3}. Ich zakres pracy obejmuje częstotliwości od \SI{24}{\mega\hertz} do \SI{1766}{\mega\hertz}. Maksymalna częstotliwość próbkowania wynosi \SI{3.2}{MS/s}, przy czym stabilną pracę uzyskuje się do około \SI{2.8}{MS/s}~\cite{sun2021}. Do odbioru ramek Mode~S wymagana jest częstotliwość próbkowania co najmniej \SI{2}{MS/s}, ponieważ impuls PPM trwa \SI{0.5}{\micro\second}~\cite{sun2021}. Odbiornik RTL-SDR Blog~V3 zawiera tuner R820T2 (oznaczany także jako R860) oraz generator kwarcowy z~kompensacją temperaturową o~stabilności \SI{1}{ppm}~\cite{rtlsdrv3www}. Próbki sygnału przekazuje do komputera przez interfejs USB w~postaci par wartości I/Q o~rozdzielczości 8~bitów~\cite{dump1090}.

Dostępne są także odbiorniki SDR o~szerszym zakresie częstotliwości, wyższej częstotliwości próbkowania i~z~możliwością nadawania, np. HackRF One i~LimeSDR~\cite{sun2021}. Wybór odbiornika dla stacji bazowej uzasadniono w~rozdziale~\ref{chap:sprzet}.

Demodulację i~dekodowanie ramek wykonuje oprogramowanie. Najbardziej znanym otwartym dekoderem Mode~S jest program dump1090, napisany w~języku C przez Salvatore Sanfilippo i~udostępniony w~2013~roku na licencji BSD~\cite{sun2021}. Program pobiera z~odbiornika próbki I/Q z~częstotliwością \SI{2}{MS/s}, oblicza ich moduł, wyszukuje preambuły ramek, demoduluje bity PPM i~sprawdza sumę kontrolną, korygując pojedyncze błędy bitowe~\cite{dump1090}. W~trybie sieciowym dump1090 udostępnia przez protokół TCP ramki zapisane szesnastkowo (port~30002) oraz zdekodowane wiadomości w~tekstowym formacie SBS, nazywanym też formatem BaseStation (port~30003)~\cite{dump1090}. Program ma wiele odmian rozwijanych niezależnie, m.in. przez firmę FlightAware~\cite{sun2021}. Przebieg odbioru sygnału przedstawiono na rysunku~\ref{rys:tor_sdr}. Dekodowanie wiadomości ADS-B, w~tym położeń zakodowanych metodą CPR, umożliwia również biblioteka pyModeS dla języka Python~\cite{sun2020pymodes}.


\begin{figure}[htbp]
  \centering
  \begin{tikzpicture}[font=\footnotesize, >=Stealth,
      blok/.style={draw, thick, fill=akcent!8, align=center, minimum height=1.1cm, text width=2.0cm, inner sep=3pt},
      wyj/.style={draw, thick, fill=black!5, align=center, minimum height=0.9cm, text width=2.3cm, inner sep=3pt}]
    \node[blok, text width=2.7cm] (ant) {antena};
    \node[blok, text width=2.7cm, right=0.65cm of ant] (tun) {tuner R820T2:\\ wzmocnienie i~przemiana\\ częstotliwości};
    \node[blok, text width=2.7cm, right=0.65cm of tun] (adc) {RTL2832U:\\ przetwornik A/C,\\ próbki I/Q};
    \draw[->, thick] (ant) -- (tun);
    \draw[->, thick] (tun) -- (adc);
    \node[blok, below=1.8cm of ant.south west, anchor=north west] (mod) {obliczenie\\ modułu próbek};
    \node[blok, right=0.65cm of mod] (pre) {wykrycie\\ preambuły};
    \node[blok, right=0.65cm of pre] (ppm) {demodulacja\\ PPM};
    \node[blok, right=0.65cm of ppm] (crc) {kontrola CRC\\ i~korekcja\\ błędów};
    \node[wyj, right=0.65cm of crc, yshift=0.7cm] (avr) {port 30002:\\ ramki (hex)};
    \node[wyj, right=0.65cm of crc, yshift=-0.7cm] (sbs) {port 30003:\\ wiadomości SBS};
    \draw[->, thick] (mod) -- (pre);
    \draw[->, thick] (pre) -- (ppm);
    \draw[->, thick] (ppm) -- (crc);
    \draw[->, thick] (crc.east) -- (avr.west);
    \draw[->, thick] (crc.east) -- (sbs.west);
    \draw[->, thick] (adc.south) -- node[right] {USB, \SI{2}{MS/s}} ++(0,-0.95) -| (mod.north);
    \begin{scope}[on background layer]
      \node[draw, dashed, black!60, fit=(ant)(tun)(adc), inner sep=7pt,
        label={[anchor=south west, black!70]north west:odbiornik RTL-SDR Blog V3}] {};
      \node[draw, dashed, black!60, fit=(mod)(pre)(ppm)(crc)(avr)(sbs), inner sep=7pt,
        label={[anchor=north west, black!70]south west:komputer: program dump1090}] {};
    \end{scope}
  \end{tikzpicture}
  \caption{Tor odbioru ramek ADS-B z~wykorzystaniem odbiornika RTL-SDR i~programu dump1090; opracowanie własne na podstawie~\cite{sun2021,dump1090}}
  \label{rys:tor_sdr}
\end{figure}

Częstotliwości \SI{1090}{\mega\hertz} odpowiada długość fali $c/f \approx \SI{27.5}{cm}$, gdzie $c$ jest prędkością światła, a~$f$ częstotliwością fali. Do odbioru sygnałów ADS-B można zatem użyć prostych anten: dipola półfalowego o~całkowitej długości około \SI{13.8}{cm} albo anteny ćwierćfalowej z~przeciwwagą, której element odbiorczy ma długość około \SI{6.9}{cm}~\cite{sun2021}.
Przy długim przewodzie antenowym sygnał o~tej częstotliwości jest silnie tłumiony. Stosuje się wtedy antenę aktywną ze wzmacniaczem niskoszumowym, zasilaną przez przewód antenowy za pośrednictwem układu \textit{bias tee}~\cite{sun2021}. Odbiornik RTL-SDR Blog~V3 ma wbudowany, włączany programowo układ \textit{bias tee} o~napięciu \SI{4.5}{V}~\cite{rtlsdrv3}.

Sygnały o~częstotliwości \SI{1090}{\mega\hertz} rozchodzą się w~zasięgu wzroku, dlatego maksymalny zasięg odbioru wynika z~krzywizny Ziemi~\cite{sun2021}:
\begin{equation}
  d = R_{\mathrm{Z}} \left( \arccos\frac{R_{\mathrm{Z}}}{R_{\mathrm{Z}} + h_{\mathrm{r}}}
    + \arccos\frac{R_{\mathrm{Z}}}{R_{\mathrm{Z}} + h_{\mathrm{t}}} \right),
  \label{eq:zasieg}
\end{equation}
\begin{gdzie}
  \item $d$ -- maksymalny zasięg odbioru, w~kilometrach,
  \item $R_{\mathrm{Z}}$ -- promień Ziemi, w~kilometrach,
  \item $h_{\mathrm{r}}$ -- wysokość anteny odbiornika nad poziomem morza, w~kilometrach,
  \item $h_{\mathrm{t}}$ -- wysokość statku powietrznego nad poziomem morza, w~kilometrach.
\end{gdzie}
Rzeczywisty zasięg jest mniejszy, ponieważ zależy także od mocy nadajnika i~kierunkowości anten~\cite{sun2021}. Wzór~\eqref{eq:zasieg} ma jednak istotną konsekwencję dla dekodowania lokalnego. Dla anteny umieszczonej na poziomie morza ($h_{\mathrm{r}} = 0$) i~$R_{\mathrm{Z}} = \SI{6371}{km}$ teoretyczny zasięg przekracza \SI{180}{NM} (\SI{333}{km}) dla statków powietrznych lecących wyżej niż około \SI{8.7}{km}, czyli na typowych wysokościach przelotowych. Dla takich statków powietrznych dekodowanie lokalne z~położeniem odbiornika jako punktem odniesienia może prowadzić do wyboru niewłaściwej strefy (podrozdział~\ref{sec:cpr}).


% ---------------------------------------------------------------------
\section{Pozycja i~czas z~odbiornika GNSS}
% =====================================================================


\label{sec:gnss}
Globalne systemy nawigacji satelitarnej (ang.~\textit{Global Navigation Satellite System}, GNSS) umożliwiają wyznaczenie położenia i~czasu w~dowolnym miejscu na Ziemi na podstawie sygnałów nadawanych przez satelity~\cite{hofmann2008}. Do systemów tych należą m.in. amerykański GPS, rosyjski GLONASS, europejski Galileo i~chiński BeiDou. Podstawowa konstelacja systemu GPS obejmuje 24~pozycje orbitalne rozmieszczone w~sześciu płaszczyznach orbit. Zapewnia to widoczność co najmniej czterech satelitów praktycznie z~każdego miejsca na Ziemi. W~praktyce utrzymuje się więcej satelitów, w~ostatnich latach około 31~\cite{gpsgov_space}.

Satelity GNSS mają na pokładzie zegary atomowe i~nadają sygnały zawierające m.in. czas nadania oraz parametry swoich orbit. Odbiornik mierzy czas propagacji sygnału każdego satelity i~wyznacza tzw. pseudoodległość, czyli odległość obarczoną błędem wynikającym z~niedokładności własnego zegara~\cite{hofmann2008}.
Dla satelity o~numerze $k$ pseudoodległość opisuje zależność~\cite{hofmann2008}
\begin{equation}
  p_k = \lVert \boldsymbol{s}_k - \boldsymbol{r} \rVert + c\,\Delta t_{\mathrm{odb}},
  \label{eq:pseudoodleglosc}
\end{equation}
\begin{gdzie}
  \item $p_k$ -- pseudoodległość do satelity $k$, w~metrach,
  \item $\boldsymbol{s}_k$ -- wektor położenia satelity $k$, w~metrach,
  \item $\boldsymbol{r}$ -- wektor położenia odbiornika, w~metrach,
  \item $c$ -- prędkość światła, w~metrach na sekundę,
  \item $\Delta t_{\mathrm{odb}}$ -- przesunięcie zegara odbiornika względem czasu systemu, w~sekundach.
\end{gdzie}
Niewiadomymi są trzy współrzędne odbiornika oraz przesunięcie jego zegara, dlatego do wyznaczenia pozycji potrzebne są pomiary z~co najmniej czterech satelitów. Ubocznym efektem obliczeń jest dokładna znajomość czasu, co pozwala wykorzystać odbiornik GNSS jako wzorzec czasu~\cite{hofmann2008}.

Odbiorniki GNSS wyrażają położenie w~układzie odniesienia WGS-84, w~którym kształt Ziemi przybliżono elipsoidą obrotową~\cite{wgs84}. Położenie opisuje się szerokością geograficzną $\varphi$, długością geograficzną $\lambda$ i~wysokością nad elipsoidą $h$. Często podaje się jednak wysokość nad poziomem morza $H$, liczoną względem geoidy, czyli powierzchni ekwipotencjalnej pola siły ciężkości zbliżonej do średniego poziomu mórz. Obie wysokości wiąże zależność~\cite{hofmann2008}
\begin{equation}
  h \approx H + N,
  \label{eq:wysokosci}
\end{equation}
\begin{gdzie}
  \item $N$ -- odstęp geoidy od elipsoidy, w~metrach.
\end{gdzie}
Odstęp ten osiąga wartości rzędu kilkudziesięciu metrów, dlatego w~obliczeniach geometrycznych wszystkie wysokości należy odnosić do tej samej powierzchni.

Odbiorniki GNSS przekazują wyniki najczęściej w~postaci tekstowych zdań standardu NMEA~0183, przesyłanych łączem szeregowym~\cite{nmea0183}. Każde zdanie rozpoczyna się znakiem \texttt{\$}, po którym następują dwuliterowy identyfikator nadawcy i~trzyliterowy typ zdania. Identyfikator \texttt{GP} oznacza dane z~systemu GPS, a~\texttt{GN} -- rozwiązanie wykorzystujące kilka systemów jednocześnie. Pola danych rozdziela się przecinkami. Zdanie kończą znak \texttt{*} i~dwucyfrowa szesnastkowa suma kontrolna, będąca wynikiem operacji XOR na kodach wszystkich znaków zawartych między \texttt{\$} a~\texttt{*}~\cite{gpsd_nmea}. Pola zdań GGA i~RMC zestawiono w~tabeli~\ref{tab:nmea}. Zdanie GGA zawiera m.in. wysokość nad poziomem morza i~odstęp geoidy od elipsoidy, co pozwala wyznaczyć wysokość nad elipsoidą ze wzoru~\eqref{eq:wysokosci}~\cite{gpsd_nmea}. Przykładowe zdania z~zaokrąglonymi współrzędnymi mają postać:
\begin{center}
  \small
  \begin{tabular}{@{}l@{}}
    \texttt{\$GPGGA,123519.00,5000.0000,N,02000.0000,E,1,08,0.9,250.0,M,40.0,M,,*64}\\
    \texttt{\$GPRMC,123519.00,A,5000.0000,N,02000.0000,E,0.0,0.0,061126,,*3B}
  \end{tabular}
\end{center}

\begin{table}[H]
  \centering
  \caption{Wybrane pola zdań GGA i~RMC standardu NMEA~0183 oraz ich wartości w~przykładowych zdaniach; opracowanie własne na podstawie~\cite{gpsd_nmea}}
  \label{tab:nmea}
  \begin{tabular}{@{}llll@{}}
    \toprule
    Zdanie & Pole & Zawartość & Przykład \\
    \midrule
    GGA & 1      & czas UTC (hhmmss.ss)                         & 12:35:19,00 \\
        & 2--3   & szerokość geograficzna (ddmm.mmmm, N/S)       & \ang{50.0} N \\
        & 4--5   & długość geograficzna (dddmm.mmmm, E/W)        & \ang{20.0} E \\
        & 6      & jakość wyznaczenia (0 -- brak, 1 -- GPS, 2 -- DGPS) & 1 \\
        & 7      & liczba wykorzystanych satelitów               & 8 \\
        & 8      & współczynnik HDOP                             & \num{0.9} \\
        & 9--10  & wysokość nad poziomem morza, w~metrach        & \num{250.0} \\
        & 11--12 & odstęp geoidy od elipsoidy, w~metrach         & \num{40.0} \\
    \midrule
    RMC & 1      & czas UTC (hhmmss.ss)                         & 12:35:19,00 \\
        & 2      & status (A -- dane poprawne, V -- niepoprawne)  & A \\
        & 3--6   & szerokość i~długość geograficzna              & jak w~GGA \\
        & 7      & prędkość względem ziemi, w~węzłach            & \num{0.0} \\
        & 8      & kąt drogi, w~stopniach                        & \num{0.0} \\
        & 9      & data (ddmmyy)                                 & 6.11.2026 \\
    \bottomrule
  \end{tabular}
\end{table}

Czas w~zdaniach NMEA jest podawany w~skali UTC, a~zdanie RMC zawiera także datę~\cite{gpsd_nmea}.
Czas systemu GPS różni się od czasu UTC o~całkowitą liczbę sekund przestępnych, którą odbiornik uwzględnia przy podawaniu czasu UTC~\cite{hofmann2008}.
Odbiornik GNSS może zatem służyć jako źródło czasu referencyjnego, niezależnego od zegara komputera. Sposób wykorzystania tego czasu do korekty znaczników czasu rejestrowanych danych opisano w~podrozdziale~\ref{sec:model}.


% ---------------------------------------------------------------------
\section{Środowisko ROS~2}
% =====================================================================
\label{sec:ros2}


Robot Operating System~2 (ROS~2) to zestaw bibliotek i~narzędzi do tworzenia oprogramowania robotów~\cite{ros2jazzy}. Mimo nazwy nie jest to system operacyjny, lecz platforma programistyczna, której podstawę stanowi warstwa pośrednicząca w~komunikacji między elementami aplikacji (ang.~\textit{middleware})~\cite{macenski2022}. ROS~2 działa w~systemie operacyjnym komputera, a~większość jego ekosystemu rozwija się wokół systemu Linux~\cite{ye2023realtime}. W~stosunku do pierwszej wersji ROS~2 zaprojektowano od nowa, m.in. z~myślą o~systemach wbudowanych, pracy w~czasie rzeczywistym i~komunikacji w~sieciach o~ograniczonej niezawodności~\cite{macenski2022}.

Oprogramowanie ROS~2 jest wydawane w~dystrybucjach. W~pracy wykorzystano dystrybucję Jazzy Jalisco, wydaną w~maju 2024~roku jako wersja o~wydłużonym okresie wsparcia, trwającym do maja 2029~roku. Najwyższy poziom wsparcia ma ona w~systemach Ubuntu~24.04, zarówno dla procesorów amd64, jak i~arm64, oraz Windows~10~\cite{ros2jazzyrelease}.

Aplikacja w~ROS~2 składa się z~węzłów (ang.~\textit{node}), czyli niezależnych jednostek programu realizujących wydzielone zadania~\cite{ros2jazzy}. Węzły wymieniają dane głównie przez tematy (ang.~\textit{topic}) w~modelu publikacji i~subskrypcji. Węzeł publikujący wysyła komunikaty (ang.~\textit{message}) do nazwanego tematu, nie wiedząc, kto je odbiera. Węzły subskrybujące dany temat otrzymują publikowane w~nim komunikaty, przy czym gwarancja ich dostarczenia zależy od ustawień jakości usługi, omówionych w~dalszej części podrozdziału~\cite{ros2jazzy}. Jeden temat może mieć wielu nadawców i~wielu odbiorców, co przedstawiono na rysunku~\ref{rys:ros2_graf}.

\begin{figure}[H]
  \centering
  \begin{tikzpicture}[font=\footnotesize, >=Stealth,
      wezel/.style={draw, thick, ellipse, fill=akcent!8, align=center, minimum height=0.95cm, inner sep=2pt},
      temat/.style={draw, thick, rectangle, fill=black!5, align=center, minimum height=0.8cm, inner sep=5pt},
      plik/.style={draw, thick, cylinder, shape border rotate=90, aspect=0.25, fill=black!5, align=center,
        minimum height=1.1cm, minimum width=1.3cm}]
    \node[wezel] (a) {węzeł A\\ (nadawca)};
    \node[temat, right=1.6cm of a] (t) {\texttt{/temat}};
    \node[wezel, right=2.0cm of t, yshift=1.25cm] (b) {węzeł B};
    \node[wezel, right=2.0cm of t] (c) {węzeł C};
    \node[wezel, right=2.0cm of t, yshift=-1.25cm] (r) {rejestrator\\ rosbag2};
    \node[plik, right=1.1cm of r] (f) {plik\\ MCAP};
    \draw[->, thick] (a) -- node[above] {publikacja} (t);
    \draw[->, thick] (t.east) -- node[sloped, above, pos=0.5] {subskrypcja} (b.west);
    \draw[->, thick] (t.east) -- (c.west);
    \draw[->, thick] (t.east) -- (r.west);
    \draw[->, thick] (r) -- node[above] {zapis} (f);
  \end{tikzpicture}
  \caption{Komunikacja węzłów ROS~2 przez temat w~modelu publikacja--subskrypcja; węzły oznaczono elipsami, a~temat -- prostokątem; rejestrator rosbag2 i~format MCAP omówiono w~dalszej części podrozdziału; opracowanie własne}
  \label{rys:ros2_graf}
\end{figure}

Komunikację między węzłami realizuje warstwa pośrednia zgodna ze standardem DDS (ang.~\textit{Data Distribution Service}). Odpowiada ona za automatyczne wykrywanie węzłów w~sieci, bez centralnego serwera, oraz za przesyłanie danych~\cite{macenski2022}. Domyślną implementacją DDS w~dystrybucji Jazzy jest eProsima Fast DDS~\cite{ros2jazzy}. Dzięki rozproszonemu wykrywaniu węzły uruchomione na różnych komputerach w~tej samej sieci mogą wymieniać komunikaty bez dodatkowej konfiguracji, o~ile mają ten sam identyfikator domeny, ustawiany zmienną środowiskową \texttt{ROS\_DOMAIN\_ID}. Przypisanie różnych identyfikatorów pozwala oddzielić od siebie grupy węzłów pracujące w~jednej sieci~\cite{ros2jazzy}.

Sposób przesyłania komunikatów określają ustawienia jakości usługi (ang.~\textit{quality of service}, QoS)~\cite{ros2jazzy}. Należą do nich m.in. niezawodność dostarczania oraz historia. W~trybie niezawodnym utracone komunikaty są wysyłane ponownie, a~w~trybie \textit{best effort} nie są. Historia określa liczbę komunikatów przechowywanych w~kolejce, nazywaną głębokością kolejki. Gdy kolejka jest pełna, najstarsze komunikaty są usuwane. Ustawienia nadawcy i~odbiorcy muszą być zgodne, w~przeciwnym razie połączenie między nimi nie powstaje~\cite{ros2jazzy}.

Opóźnienie przesyłania komunikatu między węzłami nie jest stałe. Zależy ono m.in. od implementacji DDS, rozmiaru komunikatu i~ustawień QoS. Dla małych komunikatów przesyłanych między węzłami działającymi na jednym komputerze jest ono rzędu ułamka milisekundy~\cite{ye2023realtime,kronauer2021latency}. Narzut samego ROS~2 w~stosunku do bezpośredniego użycia DDS może sięgać 50\,\% dla małych komunikatów, a~tryb niezawodny zwiększa opóźnienie o~kilkanaście procent w~porównaniu z~trybem \textit{best effort}~\cite{kronauer2021latency}.

Strukturę komunikatu opisuje plik interfejsu z~rozszerzeniem \texttt{.msg}, w~którym wymienia się pola i~ich typy~\cite{ros2jazzy}. Na tej podstawie narzędzia ROS~2 generują kod obsługujący komunikat w~językach C++ i~Python. Dostępne są standardowe typy komunikatów, np. \texttt{sensor\_msgs/NavSatFix}, przeznaczony do przekazywania pozycji z~odbiornika GNSS, oraz \texttt{sensor\_msgs/TimeReference}, przeznaczony do przekazywania czasu z~zewnętrznego źródła~\cite{ros2jazzy}. Gdy żaden typ standardowy nie odpowiada potrzebom, można zdefiniować własny typ komunikatu. Wiele typów komunikatów zawiera nagłówek ze znacznikiem czasu i~nazwą układu współrzędnych, do którego odnoszą się dane~\cite{ros2jazzy}. Znacznik czasu nadaje węzeł publikujący. Domyślnie jest to czas zegara systemowego komputera, na którym działa węzeł, a~po ustawieniu parametru \texttt{use\_sim\_time} -- czas publikowany w~temacie \texttt{/clock}~\cite{ros2jazzy}. Dokładność znaczników czasu zależy zatem od synchronizacji zegara komputera, co uwzględniono w~modelu opisanym w~podrozdziale~\ref{sec:model}.

Do rejestracji danych służy narzędzie rosbag2, które zapisuje komunikaty wybranych tematów do pliku wraz ze znacznikami czasu i~pozwala je później odtworzyć~\cite{ros2jazzy}. Odtwarzane komunikaty są publikowane w~tych samych tematach, dlatego węzły przetwarzające dane działają tak samo jak podczas pracy na żywo. Podczas odtwarzania można publikować również zegar odpowiadający czasowi nagrania. Węzły z~ustawionym parametrem \texttt{use\_sim\_time} posługują się wtedy czasem zapisu zamiast czasem bieżącym~\cite{ros2jazzy}. Od dystrybucji Iron, wydanej w~maju 2023~roku, domyślnym formatem zapisu jest MCAP, który zastąpił stosowany wcześniej format SQLite3~\cite{ros2iron}. Plik MCAP zawiera definicje zapisanych komunikatów, dzięki czemu można go odczytać także bez zainstalowanych definicji typów komunikatów~\cite{ros2iron,mcap}.


% ---------------------------------------------------------------------
\section{Istniejące rozwiązania}
% =====================================================================
\label{sec:istniejace}


Odbiór sygnałów ADS-B za pomocą tanich odbiorników SDR jest powszechny, a~budowa własnego odbiornika jest popularnym projektem hobbystycznym. Typowy odbiornik składa się z~anteny, odbiornika RTL-SDR i~dekodera, np. programu dump1090~\cite{sun2021}, uruchamianego często na komputerze jednopłytkowym Raspberry Pi~\cite{fr24build,piaware}. Odebrane statki powietrzne wyświetla się zwykle na mapie w~przeglądarce internetowej, a~zdekodowane dane są dostępne przez porty TCP programu dump1090~\cite{dump1090}. Takie odbiorniki służą przede wszystkim do obserwacji ruchu lotniczego w~okolicy. Dane nie są w~nich udostępniane jako komunikaty środowiska ROS~2, a~w~podstawowej konfiguracji nie są też zapisywane do plików.

Odbiorniki budowane przez hobbystów często przekazują dane do serwisów śledzenia lotów. Serwis Flightradar24 udostępnia instrukcję budowy odbiornika złożonego z~komputera Raspberry Pi, odbiornika ADS-B w~postaci przystawki USB oraz anteny. Na komputerze instaluje się przygotowany przez serwis obraz systemu Pi24, oparty na systemie Raspbian Lite, a~odbiornik automatycznie przesyła odebrane dane do serwisu. Według serwisu zasięg takiego odbiornika wynosi od 200 do 400~mil, czyli od około 320 do \SI{640}{km}. W~zamian za udostępnianie danych użytkownik otrzymuje bezpłatny dostęp do płatnego planu subskrypcji~\cite{fr24build}. Podobnie działa oprogramowanie PiAware, które na komputerze Raspberry Pi dekoduje wiadomości ADS-B, wyświetla je na lokalnej mapie i~przesyła do serwisu FlightAware. Jego użytkownicy otrzymują bezpłatne konto z~rozszerzonym dostępem do danych serwisu~\cite{piaware}.

Przykładem rozwiązania naukowego jest OpenSky Network, czyli sieć odbiorników naziemnych, która gromadzi surowe wiadomości Mode~S i~ADS-B w~centralnej bazie danych i~udostępnia je do celów badawczych~\cite{schafer2014}. Wszystkie opisane sieci są nastawione na gromadzenie danych z~wielu odbiorników w~jednym miejscu. Do przesyłania danych wymagają połączenia z~internetem, a~dane udostępniają przez własne serwisy.

Dane ADS-B wykorzystuje się także w~badaniach nad wizyjnym wykrywaniem statków powietrznych. Zhang i~in.~\cite{zhang2022sta} zaproponowali sieć neuronową do segmentacji statków powietrznych w~nagraniach wideo z~lotniska. Położenie odczytane z~wiadomości ADS-B służy w~niej jako informacja wstępna, która ułatwia śledzenie obiektów w~długich sekwencjach. W~zbiorze danych TartanAviation~\cite{patrikar2025tartan} zarejestrowano obrazy z~kamer zamontowanych na stałe w~obrębie lotnisk wraz z~danymi ADS-B z~własnych odbiorników. Położenia statków powietrznych interpolowano do częstotliwości rejestracji obrazu, wynoszącej \SI{24}{\hertz}, i~rzutowano na obraz za pomocą skalibrowanego modelu kamer. Rzutowane położenia posłużyły do odrzucania błędnych wykryć i~wskazania przybliżonego położenia statku powietrznego na obrazie, a~etykiety sprawdzono następnie ręcznie~\cite{patrikar2025tartan}. Przykłady te pokazują, że do automatycznego etykietowania obrazów potrzebne są zsynchronizowany czas, znane położenie odbiornika i~kamery oraz wiarygodne położenia statków powietrznych.

\begin{table}[htbp]
  \centering
  \caption{Porównanie odbiorników ADS-B o~różnym przeznaczeniu; opracowanie własne na podstawie~\cite{fr24build,piaware,schafer2014,dump1090}}
  \label{tab:rozwiazania}
  \small
  \begin{tabular}{@{}>{\raggedright\arraybackslash}p{2.8cm}>{\raggedright\arraybackslash}p{3.9cm}>{\raggedright\arraybackslash}p{3.9cm}>{\raggedright\arraybackslash}p{3.9cm}@{}}
    \toprule
    Cecha & Sieci odbiorników (Flightradar24, FlightAware, OpenSky) & Odbiornik hobbystyczny z~programem dump1090 & Stacja bazowa (niniejsza praca) \\
    \midrule
    Przeznaczenie & gromadzenie danych z~wielu odbiorników i~ich udostępnianie & obserwacja ruchu lotniczego w~okolicy odbiornika & źródło danych referencyjnych w~terenie \\
    \addlinespace
    Połączenie z~internetem & wymagane do przesyłania danych & niewymagane & niewymagane \\
    \addlinespace
    Postać udostępnianych danych & serwis internetowy lub baza danych & mapa w~przeglądarce, porty TCP & komunikaty ROS~2 w~czasie rzeczywistym \\
    \addlinespace
    Zapis danych & centralnie, w~serwisie & brak w~podstawowej konfiguracji & lokalnie, pliki rosbag2 w~formacie MCAP \\
    \bottomrule
  \end{tabular}
\end{table}

Cechy omówionych odbiorników zestawiono w~tabeli~\ref{tab:rozwiazania}. W~przeanalizowanych rozwiązaniach nie znaleziono przenośnego systemu, który w~warunkach polowych, bez dostępu do internetu, wyznaczałby położenia statków powietrznych względem własnej pozycji ustalonej przez odbiornik GNSS i~udostępniał je w~czasie rzeczywistym jako komunikaty środowiska ROS~2. Wymagania stawiane takiemu systemowi sformułowano w~podrozdziale~\ref{sec:wymagania_referencja}.
